\documentclass[tikz,border=5pt]{standalone}
\usepackage{amsmath}
\usetikzlibrary{arrows.meta,calc}

\definecolor{ink}{HTML}{1F2937}
\definecolor{slate}{HTML}{334155}
\definecolor{inputfill}{HTML}{EAF3FB}
\definecolor{inputedge}{HTML}{829DB4}
\definecolor{baselinefill}{HTML}{EEF2F6}
\definecolor{baselineedge}{HTML}{8795A5}
\definecolor{residualfill}{HTML}{FFF3DA}
\definecolor{residualedge}{HTML}{B49B63}
\definecolor{correctionfill}{HTML}{E7F3EF}
\definecolor{correctionedge}{HTML}{6B9D92}
\definecolor{correctionarrow}{HTML}{227C70}
\definecolor{evaluationfill}{HTML}{F1EFF6}
\definecolor{evaluationedge}{HTML}{9892A5}

\newcommand{\heading}[1]{%
  {\sffamily\fontsize{9.5}{11}\selectfont\bfseries #1}%
}
\newcommand{\labelfont}{\sffamily\fontsize{9}{10.5}\selectfont}

\newsavebox{\InputMath}
\newsavebox{\TargetMath}
\newsavebox{\BaselineMath}
\newsavebox{\ResidualMath}
\newsavebox{\MeanMath}
\newsavebox{\CorrectedMath}
\newsavebox{\ComparisonMath}
\newsavebox{\EvaluationMath}

\newsavebox{\InputBody}
\newsavebox{\TargetBody}
\newsavebox{\BaselineBody}
\newsavebox{\ResidualBody}
\newsavebox{\MeanBody}
\newsavebox{\CorrectedBody}
\newsavebox{\EvaluationBody}
\newsavebox{\AverageLabel}

\newlength{\ColumnA}
\newlength{\ColumnB}
\newlength{\ColumnC}
\newlength{\ColumnD}
\newlength{\AverageGap}
\newlength{\SmallGap}
\newlength{\DiagramWidth}
\newlength{\RequiredWidth}
\newlength{\UpperHeight}
\newlength{\LowerHeight}
\newlength{\TemporaryHeight}
\newlength{\RowSeparation}
\newlength{\ChannelY}
\newlength{\XA}
\newlength{\XB}
\newlength{\XC}
\newlength{\XD}

% Measure complete typeset blocks, including aligned continuation lines.
% The additional 5.2 mm supplies 2.6 mm horizontal padding on each side.
\newcommand{\columnwidth}[3]{%
  \setlength{#1}{\wd#2}%
  \ifdim\wd#3>#1
    \setlength{#1}{\wd#3}%
  \fi
  \addtolength{#1}{5.2mm}%
}
\newcommand{\includeheight}[2]{%
  \setlength{\TemporaryHeight}{%
    \dimexpr\ht#2+\dp#2+4.8mm\relax}%
  \ifdim\TemporaryHeight>#1
    \setlength{#1}{\TemporaryHeight}%
  \fi
}

\begin{document}
\color{ink}
\sffamily
\fontsize{10}{12}\selectfont

\sbox{\InputMath}{%
  $\displaystyle \mathbf{x}=(-2,-1,0,1,2)$}

\sbox{\TargetMath}{%
  $\displaystyle
  \begin{gathered}
    \mathbf{y}=2\mathbf{x}+3\cdot\mathbf{1}_5\\[0.6mm]
    \mathbf{y}=(-1,1,3,5,7)\\[0.8mm]
    \textstyle\mathbf{1}_5=(1,1,1,1,1)
  \end{gathered}$}

\sbox{\BaselineMath}{%
  $\displaystyle
  \begin{aligned}
    \widehat{\mathbf{y}}^{(0)}&=2\mathbf{x}\\[0.6mm]
                            &=(-4,-2,0,2,4)
  \end{aligned}$}

\sbox{\ResidualMath}{%
  $\displaystyle
  \begin{gathered}
    \mathbf{r}=\mathbf{y}-2\mathbf{x}\\[0.6mm]
    \mathbf{r}=(3,3,3,3,3)
  \end{gathered}$}

\sbox{\MeanMath}{%
  $\displaystyle b=\frac{1}{5}\sum_{i=1}^{5}r_i=3$}

\sbox{\CorrectedMath}{%
  $\displaystyle
  \begin{aligned}
    \widehat{\mathbf{y}}^{\,\mathrm{corr}}
      &=2\mathbf{x}+b\cdot\mathbf{1}_5\\[0.6mm]
    \widehat{\mathbf{y}}^{\,\mathrm{corr}}
      &=(-1,1,3,5,7)
  \end{aligned}$}

\sbox{\ComparisonMath}{%
  $\displaystyle
    \mathbf{y}\quad\text{vs.}\quad
    \widehat{\mathbf{y}}^{\,\mathrm{corr}}$}

\sbox{\EvaluationMath}{%
  $\displaystyle
  \begin{aligned}
    \mathrm{MSE}
      &=\frac{1}{5}\sum_{i=1}^{5}
        \left(y_i-\widehat y_i^{\,\mathrm{corr}}\right)^2\\[0.5mm]
      &=0
  \end{aligned}$}

\sbox{\InputBody}{%
  \begin{tabular}{@{}c@{}}
    \heading{Input observations}\\[1.2mm]
    \usebox{\InputMath}
  \end{tabular}}

\sbox{\TargetBody}{%
  \begin{tabular}{@{}c@{}}
    \heading{Observed targets}\\[1.2mm]
    \usebox{\TargetMath}
  \end{tabular}}

\sbox{\BaselineBody}{%
  \begin{tabular}{@{}c@{}}
    \heading{Fixed baseline}\\[1.2mm]
    \usebox{\BaselineMath}
  \end{tabular}}

\sbox{\ResidualBody}{%
  \begin{tabular}{@{}c@{}}
    \heading{Observed residuals}\\[1.2mm]
    \usebox{\ResidualMath}
  \end{tabular}}

\sbox{\MeanBody}{%
  \begin{tabular}{@{}c@{}}
    \heading{Mean residual}\\[1.2mm]
    \usebox{\MeanMath}
  \end{tabular}}

\sbox{\CorrectedBody}{%
  \begin{tabular}{@{}c@{}}
    \heading{Corrected prediction}\\[1.2mm]
    \usebox{\CorrectedMath}
  \end{tabular}}

\sbox{\EvaluationBody}{%
  \begin{tabular}{@{}c@{}}
    \heading{Evaluate on the same}\\[-0.3mm]
    \heading{five observations}\\[1.2mm]
    \usebox{\ComparisonMath}\\[1.1mm]
    \usebox{\EvaluationMath}
  \end{tabular}}

\sbox{\AverageLabel}{%
  \labelfont
  \begin{tabular}{@{}c@{}}
    average\\[-0.3mm]
    residuals
  \end{tabular}}

\columnwidth{\ColumnA}{\InputBody}{\TargetBody}
\columnwidth{\ColumnB}{\BaselineBody}{\ResidualBody}
\columnwidth{\ColumnC}{\MeanBody}{\MeanBody}
\columnwidth{\ColumnD}{\CorrectedBody}{\EvaluationBody}

% Reserve measured label width plus clearance from both neighboring nodes.
\setlength{\AverageGap}{\dimexpr\wd\AverageLabel+3.4mm\relax}

% Approximately 18 cm overall, including the standalone border.
% Expand only if the measured mathematics requires additional room.
\setlength{\DiagramWidth}{\dimexpr18cm-10pt\relax}
\setlength{\RequiredWidth}{%
  \dimexpr\ColumnA+\ColumnB+\ColumnC+\ColumnD
  +\AverageGap+11mm\relax}
\ifdim\RequiredWidth>\DiagramWidth
  \setlength{\DiagramWidth}{\RequiredWidth}
\fi
\setlength{\SmallGap}{%
  \dimexpr\DiagramWidth-\ColumnA-\ColumnB-\ColumnC-\ColumnD
  -\AverageGap\relax}
\divide\SmallGap by 2

\setlength{\XA}{\dimexpr\ColumnA/2\relax}
\setlength{\XB}{%
  \dimexpr\ColumnA+\SmallGap+\ColumnB/2\relax}
\setlength{\XC}{%
  \dimexpr\ColumnA+\SmallGap+\ColumnB+\AverageGap
  +\ColumnC/2\relax}
\setlength{\XD}{\dimexpr\DiagramWidth-\ColumnD/2\relax}

\setlength{\UpperHeight}{0pt}
\includeheight{\UpperHeight}{\InputBody}
\includeheight{\UpperHeight}{\BaselineBody}
\includeheight{\UpperHeight}{\CorrectedBody}

\setlength{\LowerHeight}{0pt}
\includeheight{\LowerHeight}{\TargetBody}
\includeheight{\LowerHeight}{\ResidualBody}
\includeheight{\LowerHeight}{\MeanBody}
\includeheight{\LowerHeight}{\EvaluationBody}

% Keep the closest upper/lower boundaries clearly separated.
\setlength{\RowSeparation}{%
  \dimexpr\UpperHeight/2+\LowerHeight/2+5.5mm\relax}
\setlength{\ChannelY}{%
  \dimexpr-\RowSeparation-\LowerHeight/2-5.8mm\relax}

\begin{tikzpicture}[
  panel/.style={
    rectangle,
    rounded corners=2mm,
    line width=0.7pt,
    inner xsep=2.6mm,
    inner ysep=2.4mm,
    outer sep=0.35pt,
    align=center,
    text=ink
  },
  supplied/.style={
    panel,fill=inputfill,draw=inputedge
  },
  fixed/.style={
    panel,fill=baselinefill,draw=baselineedge
  },
  residual/.style={
    panel,fill=residualfill,draw=residualedge
  },
  corrected/.style={
    panel,fill=correctionfill,draw=correctionedge
  },
  evaluation/.style={
    panel,fill=evaluationfill,draw=evaluationedge
  },
  dependency/.style={
    line width=0.9pt,
    -{Latex[length=1.9mm,width=1.25mm]},
    shorten <=0.35pt,
    shorten >=0.35pt,
    line cap=round,
    line join=round
  },
  slate arrow/.style={dependency,draw=slate},
  teal arrow/.style={dependency,draw=correctionarrow},
  edge label/.style={
    font=\labelfont,
    text=ink,
    fill=white,
    inner xsep=0.6mm,
    inner ysep=0.25mm,
    outer sep=0pt,
    align=center
  }
]

\node[supplied,minimum width=\ColumnA]
  (input) at (\XA,0) {\usebox{\InputBody}};

\node[supplied,minimum width=\ColumnA]
  (targets) at (\XA,-\RowSeparation) {\usebox{\TargetBody}};

\node[fixed,minimum width=\ColumnB]
  (baseline) at (\XB,0) {\usebox{\BaselineBody}};

\node[residual,minimum width=\ColumnB]
  (residuals) at (\XB,-\RowSeparation) {\usebox{\ResidualBody}};

\node[corrected,minimum width=\ColumnC]
  (mean) at (\XC,-\RowSeparation) {\usebox{\MeanBody}};

\node[corrected,minimum width=\ColumnD]
  (prediction) at (\XD,0) {\usebox{\CorrectedBody}};

\node[evaluation,minimum width=\ColumnD]
  (evaluate) at (\XD,-\RowSeparation) {\usebox{\EvaluationBody}};

% 1. Supplied inputs to the fixed baseline.
\draw[slate arrow]
  (input.east) -- (baseline.west);

% 2. Fixed baseline to observed residuals.
\draw[slate arrow]
  (baseline.south) -- (residuals.north);

% 3. Supplied targets to observed residuals.
\draw[slate arrow]
  (targets.east) -- (residuals.west);

% 4. Residuals to their scalar mean.
\draw[teal arrow]
  (residuals.east) --
  node[edge label,above=1.5mm] {\usebox{\AverageLabel}}
  (mean.west);

% 5. Baseline contribution across the open upper third column.
\draw[slate arrow]
  (baseline.east) --
  node[edge label,above=1.5mm] {$2\mathbf{x}$}
  (prediction.west);

% 6. Separate correction contribution, below the baseline dependency.
\coordinate (correction-port) at ([yshift=-5mm]prediction.west);
\coordinate (correction-elbow) at (mean.north |- correction-port);
\draw[teal arrow,rounded corners=1.1mm]
  (mean.north) -- (correction-elbow) --
  node[edge label,below=1.5mm] {$b=3$}
  (correction-port);

% 7. Corrected predictions to evaluation.
\draw[slate arrow]
  (prediction.south) -- (evaluate.north);

% 8. Independent target-data route beneath every lower-row node.
\coordinate (channel-level) at (0,\ChannelY);
\draw[slate arrow,rounded corners=1.1mm]
  (targets.south) --
  (targets.south |- channel-level) --
  node[edge label,below=1.5mm] (target-route-label)
    {observed targets: same five points}
  (evaluate.south |- channel-level) --
  (evaluate.south);

\coordinate (diagram-center) at
  (\dimexpr\DiagramWidth/2\relax,0);
\node[
  anchor=north,
  inner sep=0pt,
  outer sep=0pt,
  font=\sffamily\fontsize{9}{11}\selectfont,
  text=ink,
  align=center
] at
  ($ (diagram-center |- target-route-label.south)+(0,-2mm) $)
  {The correction is computed and evaluated on these same five observations.};

\end{tikzpicture}
\end{document}